\section{Preliminaries from algebra and complex analysis}\label{sec:prelim}
% ===============================================================

The goal of this section is a single statement, Corollary~\ref{cor:localorder}:
near any point, a nonzero meromorphic function looks like a power of a local
coordinate times a nonvanishing function, and the exponent is a well-defined
integer. That integer is what a divisor records.

\subsection{The polynomial ring and the rational function field}

The polynomial ring over $\CC$ is
\[
\CC[z] \;=\; \Big\{ \textstyle\sum_{j=0}^{n} a_j z^j \;:\; a_j \in \CC,\ n \geq 0 \Big\},
\]
an integral domain, and its field of fractions is the field of rational functions
\[
\CC(z) \;=\; \Big\{ \tfrac{P(z)}{Q(z)} \;:\; P, Q \in \CC[z],\ Q \neq 0 \Big\}.
\]
We write $\CC^\times = \CC \setminus \{0\}$ and $\CC(z)^\times$ for the group of
nonzero rational functions under multiplication.

\begin{smjproposition}[Reduced form]\label{prop:reduced}
Every $f \in \CC(z)^\times$ may be written as
\[
f(z) \;=\; c\, \frac{\prod_{i=1}^{r} (z-a_i)^{m_i}}{\prod_{j=1}^{s} (z-b_j)^{n_j}},
\qquad c \in \CC^\times,
\]
where $r, s \geq 0$ (empty products being $1$), the integers $m_i, n_j$ are
positive, and the $r+s$ points $a_1, \dots, a_r, b_1, \dots, b_s \in \CC$ are
\emph{pairwise distinct}. In such an expression the $a_i$ are precisely the zeros
of $f$ in $\CC$, with multiplicities $m_i$, and the $b_j$ are precisely its poles
in $\CC$, with orders $n_j$.
\end{smjproposition}

\begin{proof}
Write $f = P/Q$ with $P, Q \in \CC[z]$, $Q \neq 0$. Since $\CC$ is algebraically
closed, the Fundamental Theorem of Algebra factors both into linear factors,
\[
P(z) = c \prod_{i} (z - \alpha_i)^{k_i}, \qquad
Q(z) = \prod_{j} (z - \beta_j)^{\ell_j},
\]
with the $\alpha_i$ pairwise distinct and the $\beta_j$ pairwise distinct.
If some $\alpha_i$ equals some $\beta_j$, cancel the common factor
$(z-\alpha_i)^{\min(k_i,\ell_j)}$; this is legitimate because $\CC[z]$ is an
integral domain, and it strictly decreases $\deg P + \deg Q$, so after finitely
many steps no $\alpha_i$ coincides with any $\beta_j$. Discarding the factors
with exponent $0$ and renaming, we obtain the displayed form with all $r+s$
points pairwise distinct.

For the final assertion, fix a point $p \in \CC$. If $p \notin \{a_i\} \cup
\{b_j\}$, then every factor is holomorphic and nonzero at $p$, so $f(p) \in
\CC^\times$. If $p = a_i$, then $f(z) = (z-a_i)^{m_i} u(z)$ where $u$ collects
the remaining factors and the constant; $u$ is holomorphic and nonzero at $a_i$
precisely because no $b_j$ equals $a_i$. The case $p = b_j$ is identical with
$(z-b_j)^{-n_j}$ in place of $(z-a_i)^{m_i}$.
\end{proof}

\begin{remark}
The cancellation step is not a formality: without it, $f = (z-1)^2/(z-1)$
satisfies the displayed formula with $a_1 = b_1 = 1$, and the reading of the
exponents as zero and pole orders fails. The reduced form is used in an essential
way in Proposition~\ref{prop:degzero}.
\end{remark}

\begin{definition}\label{def:mult}
Let $P \in \CC[z]$ and $a \in \CC$. We say $a$ is a \emph{zero of multiplicity
$m$} of $P$ if $P(z) = (z-a)^m Q(z)$ with $Q \in \CC[z]$ and $Q(a) \neq 0$.
\end{definition}

\begin{example}\label{ex:mult}
Let $P(z) = (z-1)^2(z+2)$. At $a = 1$ we may write $P(z) = (z-1)^2 Q(z)$ with
$Q(z) = z+2$ and $Q(1) = 3 \neq 0$, so $1$ is a zero of multiplicity $2$.
At $a = -2$ we write $P(z) = (z+2) R(z)$ with $R(z) = (z-1)^2$ and $R(-2) = 9
\neq 0$, so $-2$ is a zero of multiplicity $1$. At every other point $P$ does not
vanish. The multiplicities sum to $3 = \deg P$, as
Proposition~\ref{prop:degsum} records in general.
\end{example}

\begin{remark}
Definition~\ref{def:mult} is stated for polynomials because that is all
Proposition~\ref{prop:degsum} needs. It is the special case, for $f$ a
polynomial and $m \geq 1$, of the general notion of local order established in
Corollary~\ref{cor:localorder} below, which applies to any meromorphic function
and allows $m$ to be negative.
\end{remark}

\begin{smjproposition}\label{prop:degsum}
If $P(z) = c \prod_{j=1}^{r} (z-a_j)^{m_j}$ with $c \in \CC^\times$, then
$\deg(P) = \sum_{j=1}^{r} m_j$.
\end{smjproposition}

\begin{proof}
Since $\CC[z]$ is an integral domain, the degree of a product of nonzero
polynomials is the sum of the degrees. Each factor $(z-a_j)^{m_j}$ has degree
$m_j$ and the constant $c$ has degree $0$, whence the claim.
\end{proof}

\subsection{Laurent expansions, poles, and meromorphic functions}

\begin{theorem}[Laurent expansion; {\cite[Ch.~5]{Ahlfors}}]\label{thm:laurent}
Let $A = \{z \in \CC : 0 < |z-a| < R\}$ and let $f$ be holomorphic on $A$. Then
there are unique coefficients $c_n \in \CC$ $(n \in \ZZ)$ with
\[
f(z) \;=\; \sum_{n=-\infty}^{\infty} c_n (z-a)^n \qquad (z \in A),
\]
the convergence being absolute and locally uniform on $A$. The coefficients are
given by $c_n = \frac{1}{2\pi i}\oint_{|z-a|=\rho} \frac{f(z)}{(z-a)^{n+1}}\,dz$
for any $0 < \rho < R$.
\end{theorem}

The negative-power part of the expansion classifies the singularity at $a$: it
vanishes identically for a removable singularity, has finitely many nonzero terms
for a pole, and infinitely many for an essential singularity.

\begin{definition}\label{def:pole}
Let $f$ be holomorphic on a punctured neighbourhood of $a \in \CC$. We say $a$ is
a \emph{pole of order $m$} of $f$, where $m \geq 1$, if the Laurent expansion of
$f$ at $a$ has the form
\[
f(z) = \sum_{n=-m}^{\infty} c_n (z-a)^n, \qquad c_{-m} \neq 0 .
\]
\end{definition}

Only now can we say what a meromorphic function is.

\begin{definition}\label{def:mero}
Let $U \subseteq \CC$ be open. A function $f$ is \emph{meromorphic} on $U$ if
there is a subset $S \subseteq U$, discrete and closed in $U$, such that $f$ is
holomorphic on $U \setminus S$ and every point of $S$ is a pole of $f$ in the
sense of Definition~\ref{def:pole}.
\end{definition}

\begin{remark}
The requirement that each singularity be a \emph{pole} is what the word
meromorphic contributes: essential singularities are excluded by definition.
Thus $e^{1/z}$ is not meromorphic on any neighbourhood of $0$, whereas $1/z^2$
is. Equivalently, $f$ is meromorphic on $U$ if and only if it extends to a
holomorphic map $U \to \PP$ that is not identically $\infty$; we will use this
formulation in Section~\ref{sec:sphere}.
\end{remark}

\begin{smjproposition}\label{prop:polefact}
If $a$ is a pole of order $m$ of $f$, then there is a function $u$, holomorphic
on a full neighbourhood of $a$, with
\[
f(z) = (z-a)^{-m} u(z), \qquad u(a) \neq 0 .
\]
\end{smjproposition}

\begin{proof}
By Definition~\ref{def:pole}, $f(z) = \sum_{n \geq -m} c_n (z-a)^n$ on a
punctured disc $0 < |z-a| < R$, with $c_{-m} \neq 0$. Factoring out
$(z-a)^{-m}$ gives $f(z) = (z-a)^{-m} \sum_{n \geq -m} c_n (z-a)^{n+m}$.
Reindexing by $k = n+m \geq 0$, define
\[
u(z) := \sum_{k=0}^{\infty} c_{k-m}(z-a)^{k}.
\]
This is a power series in nonnegative powers of $z-a$. It converges wherever the
original Laurent series does, hence on $0 < |z-a| < R$; but a power series
converging on a punctured disc converges on the whole disc, since its radius of
convergence is determined by its coefficients alone. Therefore $u$ is holomorphic
on $|z-a| < R$, in particular \emph{at} $a$ --- which is the point of the
reindexing, since $f$ itself was only defined on the punctured disc. Finally
$u(a) = c_{-m} \neq 0$.
\end{proof}

\subsection{Local factorisation and the local order}

\begin{smjproposition}\label{prop:zerofact}
Let $f$ be holomorphic near $a \in \CC$ and not identically zero near $a$. Then
either $f(a) \neq 0$, or there exist a unique integer $m \geq 1$ and a
holomorphic function $u$ near $a$ with
\[
f(z) = (z-a)^m u(z), \qquad u(a) \neq 0 .
\]
\end{smjproposition}

\begin{proof}
Expand $f(z) = \sum_{n \geq 0} c_n (z-a)^n$ near $a$. If $c_0 = f(a) \neq 0$
there is nothing to prove. Otherwise, since $f \not\equiv 0$, not all $c_n$
vanish; let $m \geq 1$ be least with $c_m \neq 0$. Then
$f(z) = (z-a)^m \sum_{k \geq 0} c_{k+m}(z-a)^{k}$, and as in
Proposition~\ref{prop:polefact} the series $u(z) := \sum_{k \geq 0}
c_{k+m}(z-a)^k$ is holomorphic near $a$ with $u(a) = c_m \neq 0$.

Uniqueness is proved in Lemma~\ref{lem:unique} below, in the generality we shall
need.
\end{proof}

\begin{smjlemma}[Uniqueness of the local exponent]\label{lem:unique}
Let $a \in \CC$ and suppose
\[
(z-a)^{m} u(z) \;=\; (z-a)^{n} v(z)
\]
on a punctured neighbourhood of $a$, where $m, n \in \ZZ$ and $u, v$ are
holomorphic and nonvanishing at $a$. Then $m = n$ and $u = v$.
\end{smjlemma}

\begin{proof}
Suppose $m \neq n$; without loss of generality $m < n$. Dividing by
$(z-a)^{m}$ --- legitimate for $z \neq a$ --- gives $u(z) = (z-a)^{n-m} v(z)$ on
a punctured neighbourhood. Both sides are holomorphic at $a$ and agree off $a$,
so by continuity they agree at $a$ as well; since $n - m \geq 1$ the right-hand
side vanishes at $a$, giving $u(a) = 0$, contrary to hypothesis. Hence $m = n$,
and then $u = v$ off $a$ and so, again by continuity, everywhere near $a$.
\end{proof}

\begin{smjcorollary}[Local order]\label{cor:localorder}
Let $f$ be meromorphic and not identically zero near $a \in \CC$. Then there are
a unique $m \in \ZZ$ and a function $u$, holomorphic and nonvanishing at $a$,
with $f(z) = (z-a)^m u(z)$ near $a$. We call $m$ the \emph{order} of $f$ at $a$.
\end{smjcorollary}

\begin{proof}
Existence: if $f$ has a pole at $a$ apply Proposition~\ref{prop:polefact} (with
$m < 0$); otherwise $f$ is holomorphic at $a$ and we apply
Proposition~\ref{prop:zerofact}, taking $m = 0$ and $u = f$ when $f(a) \neq 0$.
Uniqueness is Lemma~\ref{lem:unique}.
\end{proof}

Thus zeros and poles alike are measured by a single integer, positive for zeros,
negative for poles, zero when the function is finite and nonvanishing. This is
the fact that makes the divisor construction possible.

% ===============================================================
